\documentclass[11pt,a4paper]{article}
\usepackage[margin=25mm]{geometry}
\usepackage{amsmath,amssymb,amsthm,hyperref}
\newtheorem{lemma}{Lemma}
\newtheorem{theorem}[lemma]{Theorem}
\newtheorem{corollary}[lemma]{Corollary}
\newcommand{\PSL}{\operatorname{PSL}}
\newcommand{\PGL}{\operatorname{PGL}}
\newcommand{\Aut}{\operatorname{Aut}}
\newcommand{\gen}[1]{\langle#1\rangle}
\title{Cubic field automorphisms and Wilson trialities}
\author{Human Compression project: L023 working proof}
\date{}
\begin{document}\maketitle

This is an unreviewed replacement for the odd-characteristic generation
argument of Leemans--Stokes, arXiv:2208.08215v2, Theorem~3.
It does not yet recover that paper's polynomial table or its assertions
about the rank-four incidence construction. No complete-original,
conjecture, formalization or compression-factor credit is assigned.
The source already states the existence result proved below; our
contribution is the explicit bound on all obstructing proper subfields.

\section{Structural inputs and a subfield count}

We retain Dickson's subgroup classification for $\PSL_2$ over a finite
field of odd characteristic, in the following precise form. A proper
subgroup is contained in a point stabilizer or is cyclic, dihedral,
isomorphic to $A_4,S_4,A_5$, or is a conjugate of a subfield group
$\PSL_2(p^d)$ or $\PGL_2(p^d)$. In the last case the defining degree
$d$ divides the ambient field degree. The $\PGL$ case belongs to
the ambient $\PSL$ only when the extension degree is even.
We also retain the standard automorphism description
$\Aut(\PSL_2(p^m))=\PGL_2(p^m)\rtimes\gen F$, with $F$ coefficient
squaring replaced here by coefficient raising to $p$.
The ambient parameters used below satisfy $p^m\geq27$, so the small
isomorphisms introduce no exception to this description.

The elementary projective-line facts used are these: the involutions
of $\PSL_2(r)$ have one $\PGL_2(r)$-conjugacy class; its Klein four
subgroups also form one such class and have normalizer $S_4$ in
$\PGL_2(r)$; an involution centralizer in $\PGL_2(r)$ is a dihedral
torus normalizer. These familiar consequences of the same classification
are retained inputs. The centralizer of a nonidentity element of odd
order is cyclic, except for a unipotent element, whose centralizer is
an elementary abelian group of odd order. We will only use the resulting
fact that such a centralizer contains at most one involution.

Write $M_d$ for a conjugate of the canonical $\PGL_2(p^d)$ inside
$\PGL_2(p^m)$, where $d\mid m$.

\begin{lemma}[Subfields through a Klein four subgroup]\label{four}
Fix $V\cong C_2^2$ in $\PSL_2(p^m)$. There are at most four
conjugates $M_d$ containing $V$. If $m/d$ is odd, there is exactly
one. In particular, if $m=3n$, every such $M_d$ with $d\mid n$
lies in the unique conjugate $M_n$ containing $V$.
\end{lemma}
\begin{proof}
First $N_{\PGL_2(p^m)}(M_d)=M_d$. The ambient fixed points of
the nonidentity unipotent elements of $M_d$ are exactly its projective
subline. A normalizer preserves that subline. Use the sharply
three-transitive action of $M_d$ on it to make a normalizing projectivity
fix three points; it is then the identity. This proves the assertion.

There are two types of Klein four subgroup in $\PGL_2(r)$.
Those contained in $\PSL_2(r)$ form the retained single class with
normalizer of order $24$. For clarity, the other type can be read
directly in an involution centralizer. The determinant square-class
map on such a subgroup has kernel $\gen z$ of order two, with
$z\in\PSL_2(r)$. In $C_{\PGL_2(r)}(z)=\gen t\rtimes\gen j$,
the torus has order $r-\varepsilon$, divisible by four, $j$ inverts
$t$, and $z=t^{(r-\varepsilon)/2}$. Both involutions outside
$\gen z$ are reflections. A torus generator has nonsquare determinant,
so the reflections outside $\PSL_2(r)$ constitute one of the two
reflection classes of this dihedral group. They give one conjugacy
class of Klein four subgroups not contained in $\PSL_2(r)$.
The normalizer of one has order eight: precisely four torus elements
and their four products with one reflection preserve its reflection
pair. A normalizer must fix the distinguished element $z$, so this
calculation inside its centralizer is the full normalizer calculation.

Let $X=\PGL_2(p^m)$. Count pairs $(M_d,W)$ with $W\leq M_d$
an $X$-conjugate of $V$. For a fixed conjugacy class of Klein four
subgroups in $M_d$, with normalizer order $e$, the number of
conjugates $M_d$ containing $V$ contributed by that class is
\[
 \frac{|X|}{|M_d|}\frac{|M_d|}{e}
 \bigg/\frac{|X|}{24}=\frac{24}{e}.
\]
The class inside its own $\PSL_2(p^d)$ contributes one. The other
class either contributes three or contributes nothing, according as
it lies in the ambient $\PSL_2(p^m)$ or not. Under an odd extension,
determinant square classes are preserved, so it contributes nothing.
This gives the stated bounds.

If $d\mid n$, extend any conjugate $M_d$ to the correspondingly
conjugated $\PGL_2(p^n)$. It still contains $V$. Since $m/n=3$
is odd, this larger group is the unique $M_n$ just proved.
\end{proof}

We use the finite-field form of Hilbert's theorem~90 for matrices:
a semilinear map on a finite-dimensional vector space over
$\mathbb F_{r^3}$ with field part a generator of
$\operatorname{Gal}(\mathbb F_{r^3}/\mathbb F_r)$ and cube the
identity is conjugate to the coordinate field map. This is a retained
descent theorem, not a new classification premise hidden in the counts.
The analogous assertion for projective semilinear transformations follows
from it. Indeed the product of a matrix lift with its two twisted copies
is a scalar in $\mathbb F_r^*$; multiply the lift by a scalar of the
inverse norm, which exists because the finite-field norm is surjective.
Matrix descent then applies. Consequently the fixed subgroup in
$\PGL_2(r^3)$ is a conjugate of $\PGL_2(r)$.

\begin{lemma}[Involution counts]\label{counts}
For odd $r$, the number of involutions in $\PGL_2(r)$ is $r^2$.
The number in $\PSL_2(r)$ is $r(r+\varepsilon)/2$, where
$\varepsilon=1$ for $r\equiv1\pmod4$ and $\varepsilon=-1$
for $r\equiv3\pmod4$.
\end{lemma}
\begin{proof}
The split and nonsplit involution centralizers in $\PGL_2(r)$
have respective orders $2(r-1)$ and $2(r+1)$. Their class sizes
are $r(r+1)/2$ and $r(r-1)/2$, whose sum is $r^2$.
Exactly the first class has square determinant when $r\equiv1
\pmod4$, and exactly the second when $r\equiv3\pmod4$.
This gives the $\PSL$ count as well.
\end{proof}

\section{Excluding every proper generated subgroup}

Let $p$ be odd, $q=p^n$, $G=\PSL_2(q^3)$ and
$X=\PGL_2(q^3)$. Choose a Klein four subgroup
$V\leq\PSL_2(p)$ and an element $s\in\PSL_2(p)$ of order
three which cycles its three involutions. Such an element is supplied
by the $A_4$ in the normalizer of $V$, also when $p=3$.
Put
\[
                   \alpha=\operatorname{conj}(s)F^n.
\]
The two factors commute and $\alpha$ has order three. Its image
in the field-automorphism quotient is nonidentity. Matrix descent
shows that $H=C_G(\alpha)$ is a conjugate of $\PSL_2(q)$:
an order-three lift of $s$ in $\operatorname{SL}_2(p)$ has norm
$s^3=1$, and the cubic extension preserves determinant square classes.

\begin{lemma}\label{bad}
For any involution $u\in H$, a proper subgroup
$J=\gen{V,u}<G$ is contained in an $\alpha$-stable conjugate
of some $M_d=\PGL_2(p^d)$ with $d\mid3n$ and $d<3n$.
\end{lemma}
\begin{proof}
The group $J$ is $\alpha$-stable, and $\alpha$ cycles $V$ while
fixing $u$. It cannot be in a point stabilizer: in odd characteristic
the latter's Sylow two-subgroups are cyclic, so cannot contain $V$.
It cannot be cyclic. A dihedral group containing $V$ has an even
rotation subgroup, and its central involution is one of the three
involutions of $V$. Unless the whole group is $V$, that central
involution is characteristic and cannot be cycled with the others.
If $J=V$, it has no involution fixed by $\alpha$.

For $J\cong A_4$, an automorphism of order three has no fixed
involution, as it cycles the three involutions. For $J\cong S_4$
or $A_5$, an automorphism of order three is inner, induced by a
three-cycle, whose centralizer in $J$ has no involution. These
cases also contradict the fixed element $u$.
The retained subgroup classification therefore makes $J$ an actual
subfield group $\PSL_2(p^d)$ or $\PGL_2(p^d)$.
The ambient normalizer of the former is its corresponding $M_d$,
and the latter normalizes itself. The same subline argument as in
Lemma~\ref{four} proves these normalizer statements; for the former,
the normalizer contains $M_d$, which can be used to adjust three points.
Since $\alpha$ normalizes $J$, it normalizes this $M_d$ as well.
\end{proof}

\begin{theorem}[Generating fixed involutions]\label{generation}
There is an involution $u\in C_G(\alpha)$ such that
$\gen{V,u}=G$. More precisely, if $n\geq2$, the number of
such involutions is at least
\[
 \frac{q(q+\varepsilon)}2-1
       -4\sum_{v=1}^{\lfloor n/2\rfloor}p^{2v}>0,
 \qquad \varepsilon=(-1)^{(q-1)/2}.
\]
For $n=1$ at most one involution of $C_G(\alpha)$ fails to
generate together with $V$.
\end{theorem}
\begin{proof}
Call an involution of $H$ bad if it generates a proper group with
$V$. By Lemma~\ref{bad}, each bad involution lies in a normalized
proper $M_d$ containing $V$.

All degrees $d\mid n$ can be treated together. By
Lemma~\ref{four}, their groups lie in the unique $M_n$ containing
$V$, namely the standard $\PGL_2(q)$. On this group $F^n$ is
trivial, so $\alpha$ acts as conjugation by $s$. Its fixed subgroup
contains at most one involution, by the retained odd-element
centralizer fact. Thus all these degrees together contribute at most
one bad involution, with no repeated charge for nested subfields.

If $d\mid3n$ but $d\nmid n$, its three-adic valuation is exactly
one larger than that of $n$. Thus
\[
 d=3v,\qquad \gcd(d,n)=v,\qquad v\mid n.
\]
For a proper degree $d<3n$ one has $n/v\geq2$, hence
$v\leq n/2$. The restriction of $F^n$ to the degree-$d$ field
has order three. On an $\alpha$-stable conjugate $M_d$ the
restriction of $\alpha$ has this same nontrivial field component
and order three. Projective matrix descent identifies its fixed
subgroup with $\PGL_2(p^v)$, so at most $p^{2v}$ of its elements
are involutions. There are at most four such conjugates containing
$V$, by Lemma~\ref{four}.

Summing over all integers $1\leq v\leq\lfloor n/2\rfloor$
only enlarges the union of bad candidates. This proves the bound.
For $n=1$ the sum is empty and $|\operatorname{Inv}(H)|\geq3$.
For $n\geq2$ put $b=\lfloor n/2\rfloor$. Then
\[
 \sum_{v=1}^{b}p^{2v}<\frac{p^2}{p^2-1}p^{2b}
 \leq\frac98q.
\]
Except for $q=9$, one has $q\geq25$, and therefore
$q(q-1)/2>1+(9/2)q$. For $q=9$ the exact bound is
$45-1-4\cdot9=8>0$. This proves positivity in every case.
\end{proof}

\section{The Wilson class}

We retain the usual correspondence between quotients of
$\gen{r_0,r_1,r_2\mid r_i^2=(r_0r_2)^2=1}$ and reflexible
maps, together with Wilson's permutation action on
$r_0,r_2,r_0r_2$. This is a map correspondence; no intersection
property for a polytopal poset is being inferred from generation.

\begin{corollary}
For every odd prime power $q$, the group $\PSL_2(q^3)$ is
the automorphism group of a reflexible map of Wilson class~III.
\end{corollary}
\begin{proof}
Take distinct $\rho_0,\rho_2\in V\setminus1$ in the cyclic
order of $\alpha$, and put $\rho_1=u$ from
Theorem~\ref{generation}. They are involutions, the endpoints
commute, and they generate $G$. The automorphism $\alpha$
fixes $\rho_1$ and cycles the three nonidentity elements of $V$,
so the resulting map has a Wilson triality.

Suppose it also had a Wilson duality. Together with the triality,
the corresponding automorphisms would give an $S_3$ acting on
$V\setminus1$ and fixing $u$. This action is faithful: an
automorphism fixing $V$ pointwise and $u$ fixes generators of $G$.
But $\operatorname{Out}(G)\cong C_2\times C_{3n}$ is abelian.
The image of the commutator subgroup $A_3$ of this $S_3$ must
be trivial there, whereas its generator $\alpha$ has nontrivial
field component. This contradiction excludes all Wilson dualities.
Thus the map belongs to class~III.
\end{proof}

\paragraph{Source boundary.}
The source's Theorem~3 for even $q$ explicitly retains the earlier
Jones--Poulton theorem. The proof above supplies all odd $q$ without
claiming that every proper $\alpha$-stable generated subgroup must
lie in the degree-$n$ subfield. Degrees $d\nmid n$ are instead
counted and excluded collectively. Its scope is map existence and
the absence of Wilson dualities; the source's polynomial solution
counts, rank-four flag completion, chamber-orbit counts, and full
geometric correlation assertions remain separate obligations.
The source authors retain attribution for the stated existence
theorem. The finite-field descent and subgroup classification are
explicitly borrowed; the normalized-subfield union bound is the
present replacement argument. No finite enumeration is a premise.
\end{document}
